\documentclass[11pt]{article}

\usepackage[margin=1in]{geometry}
\usepackage{amsmath,amssymb,amsthm}
\usepackage{array}

\newtheorem{theorem}{Theorem}[section]
\newtheorem{lemma}[theorem]{Lemma}
\newtheorem{proposition}[theorem]{Proposition}
\newtheorem{corollary}[theorem]{Corollary}
\newtheorem{assumption}[theorem]{Assumption}
\theoremstyle{definition}
\newtheorem{definition}[theorem]{Definition}
\theoremstyle{remark}
\newtheorem{remark}[theorem]{Remark}

\newcommand{\area}{\operatorname{area}}
\newcommand{\dinv}{\operatorname{dinv}}
\newcommand{\defc}{\operatorname{defc}}
\newcommand{\UP}{\operatorname{UP}}
\newcommand{\DOWN}{\operatorname{DOWN}}
\newcommand{\East}{\operatorname{East}}
\newcommand{\West}{\operatorname{West}}
\newcommand{\inj}{\operatorname{inj}}
\newcommand{\bk}{\operatorname{bk}}
\newcommand{\Cat}{\operatorname{Cat}}
\newcommand{\NN}{\mathbb N}

\title{The Modulo-One String Decomposition}
\author{}
\date{}

\begin{document}
\maketitle

\section{Statement and conventions}

Fix integers
\[
  \tau\ge2,
  \qquad s\ge6,
  \qquad r=\tau s+1,
  \qquad M_s=\tau\binom{s}{2}.
\]
A \emph{normalized step-$\tau$ Dyck word} of length $s$ is a word
\[
  D=(x_0,\ldots,x_{s-1})
\]
of nonnegative integers such that
\[
  x_0=0,
  \qquad x_{i+1}\le x_i+\tau.
\]
Its statistics are
\[
  \area(D)=\sum_i x_i
\]
and
\[
  \dinv_\tau(D)=\sum_{i<j}d_\tau(x_i,x_j),
\]
where
\[
 d_\tau(u,v)=
 \begin{cases}
   (u+\tau-v)_+,&u\le v,\\
   (v+\tau+1-u)_+,&u>v.
 \end{cases}
\]
Here $z_+=\max(0,z)$.  We put
\[
  \defc_\tau(D)=M_s-\area(D)-\dinv_\tau(D).
\]
For a fixed defect $d$, write
\[
  L_d=\left\lfloor\frac{M_s-d}{2}\right\rfloor.
\]
Thus $D$ is in the weak lower half of its defect layer precisely when
$\area(D)\le L_d$, and it is strictly below the middle when
$\area(D)\le L_d-1$.

We write $A:B$ for concatenation.  An occurrence $x_j=e>0$ is
\emph{extractable} if exactly one earlier occurrence lies in
\[
  [\max(0,e-\tau),e)
\]
and deleting $x_j$ leaves a normalized word.  If $0<j<s-1$, the latter
condition is the splice inequality
\[
  x_{j+1}\le x_{j-1}+\tau.
\]
The extraction convention always chooses the leftmost extractable
occurrence.

Conversely, if $e>0$ and a normalized word contains an occurrence in
\[
  [\max(0,e-\tau),e),
\]
then $\inj_e$ inserts $e$ immediately after the first such occurrence.
This is the inverse of extracting the inserted occurrence whenever the
resulting insertion is admissible.

A \emph{full skeleton} is a normalized word with no extractable occurrence.
For $s\ge4$, let
\[
  \epsilon_{s,\tau}
  =(0,0,1,\underbrace{0,\ldots,0}_{s-4\text{ entries}},\tau).
\]
A \emph{special skeleton} is a full skeleton other than
$\epsilon_{s,\tau}$.  The exceptional skeleton and the word
$(0,\ldots,0,\tau)$ both have defect $\tau(s-2)$, one beyond the range
treated below.

The rational $q,t$-Catalan polynomial in this normalization is
\[
  \Cat_{\tau s+1,s}(q,t)
  =\sum_D q^{\dinv_\tau(D)}t^{\area(D)},
\]
where the sum is over normalized step-$\tau$ Dyck words of length $s$.

\subsection{The local moves}

On a three-window $(a,c,p)$, the map $\East_3$ is defined when
\[
  c\le p+\tau
\]
and is then the identity.  Suppose instead that $c>p+\tau$ in a
five-window $(a,b,c,p,q)$.  Put
\[
  \bk_\tau(b,c)=
  \begin{cases}
    (c,b),&b>c+\tau,\\
    (b,c),&b\le c+\tau.
  \end{cases}
\]
The map $\East_5$ is defined as follows.
\begin{enumerate}
\item If $b\le p+\tau$, it is defined when $b\le q+\tau$ and sends
\[
  (a,b,c,p,q)\longmapsto(a,p,c,b,q).
\]
\item If $b>p+\tau$, write $(b',c')=\bk_\tau(b,c)$.  It is defined when
$c'\le q+\tau$ and sends
\[
  (a,b,c,p,q)\longmapsto(a,p,b',c',q).
\]
\end{enumerate}
The maps $\West_3$ and $\West_5$ are obtained by reversing the input,
applying the corresponding East map, and reversing the output.

\subsection{The three assumptions}

The proof below isolates the three facts being assumed.

\begin{assumption}[Special-skeleton position]
\label{ass:skeleton-area}
Every special skeleton $S$ satisfies
\[
  \area(S)\le\defc_\tau(S).
\]
\end{assumption}

\begin{assumption}[East5 nonfailure]
\label{ass:east5}
Every $\East_5$ stage arising from a strictly lower-half word of defect at
most $\tau(s-2)-1$ is in the domain of $\East_5$.
\end{assumption}

\begin{assumption}[West5 nonfailure]
\label{ass:west5}
Every $\West_5$ stage arising from a weakly lower-half nonskeleton of defect
at most $\tau(s-2)-1$ is in the domain of $\West_5$.
\end{assumption}

We also assume the known symmetry
\[
  \Cat_{\tau s+1,s}(q,t)=\Cat_{\tau s+1,s}(t,q).
\]

\begin{theorem}[Modulo-one skeleton decomposition]
\label{thm:main}
Under the preceding assumptions, put
\[
  B=\tau(s-2)-1.
\]
Then
\begin{align*}
 &\sum_{\substack{D\text{ normalized}\\\defc_\tau(D)\le B}}
 q^{\dinv_\tau(D)}t^{\area(D)}\\
 &\qquad=
 \sum_{\substack{S\text{ special skeleton}\\\defc_\tau(S)\le B}}
 \frac{
 q^{\dinv_\tau(S)+1}t^{\area(S)}
 -q^{\area(S)}t^{\dinv_\tau(S)+1}
 }{q-t}.
\end{align*}
\end{theorem}

\section{Common structural lemmas}

Define the complementary kernel
\[
 K_\tau(u,v)=\tau-d_\tau(u,v)
 =
 \begin{cases}
   \min(\tau,v-u),&u\le v,\\
   \min(\tau,u-v-1),&u>v.
 \end{cases}
\]
For $X=(x_0,\ldots,x_{N-1})$, put
\[
  \lambda_j(X)=\sum_{i<j}K_\tau(x_i,x_j)-x_j.
\]
Then
\begin{equation}
  \defc_\tau(X)=\sum_{j=0}^{N-1}\lambda_j(X),
  \label{eq:defect-lambda}
\end{equation}
where the defect on the left is computed using
$\tau\binom N2$.

\begin{lemma}[Record payment]
\label{lem:record-payment}
Let
\[
  0=r_0<r_1<\cdots<r_h=m,
  \qquad r_i-r_{i-1}\le\tau,
\]
be a chain of record values.  For $0\le u\le m+\tau$,
\[
  \sum_{r_i\le u}K_\tau(r_i,u)\ge u.
\]
Consequently every lambda increment of a normalized word is nonnegative.
If a normalized word begins with two zeros and has record peak $m$, then
\[
  \defc_\tau(X)\ge m.
\]
\end{lemma}

\begin{proof}
Let $r_j$ be the largest record at most $u$.  For $i<j$,
\[
  K_\tau(r_i,u)\ge r_{i+1}-r_i,
\]
while $K_\tau(r_j,u)=u-r_j$.  The sum telescopes to $u$.  In a normalized
word, the preceding records at or below the current entry therefore pay the
subtraction in its lambda increment.

If the word begins with two zeros, the second zero is not needed for these
payments.  At a positive record $r_i$ it contributes
$K_\tau(0,r_i)=\min(\tau,r_i)$.  Thus the unused contributions of the second
zero are at least
\[
  \sum_{i=1}^h\min(\tau,r_i)
  \ge\sum_{i=1}^h(r_i-r_{i-1})=m.
\]
Equation~\eqref{eq:defect-lambda} completes the proof.
\end{proof}

\begin{lemma}[Failed-splice propagation]
\label{lem:propagation}
Let $X=(x_0,\ldots,x_{N-1})$ be normalized and suppose $x_1>0$.  If none
of $x_1,\ldots,x_h$ is extractable, then the entries through $x_{h+1}$ are
strictly increasing records and
\begin{equation}
  x_{i+2}>x_i+\tau
  \qquad(0\le i<h).
  \label{eq:rapid}
\end{equation}
In particular:
\begin{enumerate}
\item every full skeleton of length at least two begins with $(0,0)$;
\item a leftmost extractable occurrence is a record;
\item if appending an entry validly blocks an extractable final occurrence,
the unique-predecessor chain continues to the appended entry, which is
extractable at the new right boundary.
\end{enumerate}
\end{lemma}

\begin{proof}
The occurrence $x_1$ has the unique predecessor $x_0=0$.  If it is not
extractable, its splice fails, so $x_2>x_0+\tau$.  Normalization gives
$x_1\le x_0+\tau$ and $x_2\le x_1+\tau$, whence $x_2>x_1$ and $x_1$ is
the unique predecessor of $x_2$.  The same argument iterates.  At the right
boundary a unique-predecessor occurrence is automatically extractable, which
proves the three consequences.
\end{proof}

\begin{lemma}[Extraction transport]
\label{lem:transport}
Let $e$ be the leftmost extractable occurrence of a normalized word $D$,
let $C$ be obtained by deleting it, and put $p=e-1$.  Then
\begin{align}
 \area(C:p)&=\area(D)-1,\notag\\
 \dinv_\tau(C:p)&=\dinv_\tau(D)+1,\label{eq:transport}\\
 \defc_\tau(C:p)&=\defc_\tau(D).\notag
\end{align}
These identities remain valid even if the final adjacency of $C:p$ fails.
\end{lemma}

\begin{proof}
Write $D=A:e:B$.  For $b\in B$,
\[
  d_\tau(e,b)=d_\tau(b,e-1).
\]
For $a\in A$, the difference
\[
  d_\tau(a,e-1)-d_\tau(a,e)
\]
is $1$ precisely when $a$ lies in the predecessor interval of $e$, is
$-1$ when $e<a\le e+\tau$, and is zero otherwise.  By
Lemma~\ref{lem:propagation}, $e$ is a record, so the negative case does not
occur.  Extractability gives exactly one predecessor and hence a total dinv
increase of $1$.  The area decreases by $1$, proving all three identities.
\end{proof}

\begin{lemma}[Double-zero boundary ledger]
\label{lem:boundary-ledger}
The following four statements hold.
\begin{enumerate}
\item If a normalized length-$s$ word begins with $(0,0)$ and only its final
occurrence is extractable, then its defect is at least $\tau(s-2)$.
\item Let $R$ be a full normalized word of length $s-1$, and let $p$ be the
lowered value arising at a first-extraction stage.  If $R:p$ is invalid or is
not full, then
\[
  \defc_\tau(R:p)\ge\tau(s-2).
\]
\item Let $C:p$ be a first-extraction stage, where $|C|=s-1$, $C$ begins
with $(0,0)$, $C[-1]>p+\tau$, and $C$ has no extractable occurrence outside
its final two positions.  Then
\[
  \defc_\tau(C:p)\ge\tau(s-2).
\]
\item Let $R:p:q$ be a second-extraction stage, where $|R|=s-2$, $R$ begins
with $(0,0)$, $p>R[-1]+\tau$, and $R$ has no extractable occurrence outside
its final position.  Then
\[
  \defc_\tau(R:p:q)\ge\tau(s-2).
\]
\end{enumerate}
\end{lemma}

\begin{proof}
All four assertions are the same kernel ledger.  At every occurrence use the
preceding record chain to pay the negative part of its lambda increment, as
in Lemma~\ref{lem:record-payment}.  The second initial zero is not used in
those payments.  Scan the remaining occurrences from left to right.  If an
occurrence has two or more predecessors, the corresponding two kernel
triangles provide a full $\tau$-unit.  If it has a unique predecessor and is
not extractable, Lemma~\ref{lem:propagation} forces its outgoing splice to
fail.  The unpaid portion of the same $\tau$-unit is carried to the next
record.  Record jumps are at most $\tau$, so these carries are nonnegative
and telescope.

At the right boundary the last carry is closed, respectively, by the
extractable final entry, by the invalid or fullness-obstructing append, by
the far pair $C[-1]>p+\tau$, or by the reverse far pair
$p>R[-1]+\tau$ together with the two extraction anchors.  The ledger contains
$s-2$ full $\tau$-units.  In the two local-window forms the count is
\[
  \underbrace{\tau}_{\text{second zero}}
  +\underbrace{(s-5)\tau}_{\text{interior}}
  +\underbrace{2\tau}_{\text{terminal block}}
  =\tau(s-2).
\]
The first two forms merely move one entry between the interior and terminal
blocks.  All omitted kernel terms are nonnegative.

For the first assertion one may see the count explicitly as follows.  If
$e$ is the final extractable value, then $e\ge\tau+1$.  Pair every nonrecord
middle occurrence $u$ with $0$ and $e$ and use
\[
  K_\tau(0,u)+K_\tau(u,e)\ge\tau.
\]
The record chain supplies the same $\tau$-unit for each record, and the
second zero supplies the final unit.  There are $s-2$ units in total.
\end{proof}

\begin{lemma}[Rapid-chain estimates]
\label{lem:rapid}
Let
\[
  R=(r_0,\ldots,r_{m-1}),
  \qquad r_0=0,
\]
be strictly increasing and normalized, and suppose
\[
  r_{i+2}>r_i+\tau.
\]
If $c=r_{m-1}$, then
\begin{align}
 \dinv_\tau(R)&=\tau(m-1)-c,\label{eq:rapid-dinv}\\
 \sum_i d_\tau(r_i,u)&\le2\tau
 \quad\text{for every }u\ge0.\label{eq:rapid-target}
\end{align}
If $m\ge5$, then
\begin{equation}
  \area(R)>\tau(m-1)\ge\dinv_\tau(R).
  \label{eq:rapid-area}
\end{equation}
If $m\ge3$, then
\begin{equation}
  \area(R)+3c+2>\tau(m+1).
  \label{eq:short-rapid}
\end{equation}
\end{lemma}

\begin{proof}
Only adjacent pairs of $R$ contribute to dinv, and their contributions
 telescope:
\[
 \sum_{i=0}^{m-2}\bigl(\tau-(r_{i+1}-r_i)\bigr)
 =\tau(m-1)-c.
\]
For a fixed target $u$, split the contributing records at $u$.  On each side
their contributions form a truncated triangle of slope one.  The rapid
condition prevents three records on the same side from occupying an interval
of width $\tau$; pairing the two closest possible records on the two sides
gives the bound $2\tau$.

Integrality gives
\[
  r_{2k}\ge k(\tau+1),
  \qquad r_{2k+1}\ge1+k(\tau+1),
\]
which implies~\eqref{eq:rapid-area}.  For~\eqref{eq:short-rapid}, the minimal
length-three chain $(0,1,\tau+1)$ has surplus $7$.  If a chain ending at $c$
is extended by $c'>c$, the surplus in~\eqref{eq:short-rapid} increases by
\[
  4c'-3c-\tau\ge c+4-\tau\ge5.
\]
Induction completes the proof.
\end{proof}

\section{The three UP lemmas}

Throughout this section, let $D=(x_0,\ldots,x_{s-1})$ be normalized and
suppose
\begin{equation}
  0\le d=\defc_\tau(D)\le\tau(s-2)-1,
  \qquad \area(D)\le L_d-1.
  \label{eq:up-hyp}
\end{equation}
Then
\begin{equation}
  \dinv_\tau(D)-\area(D)\ge2.
  \label{eq:strict-lower-gap}
\end{equation}

\begin{lemma}[Validity of the full-skeleton branch]
\label{lem:up-skeleton}
Suppose $D$ is full.  Then $x_{s-1}+1$ has an injection anchor in the prefix.
Consequently the full-skeleton branch of $\UP$ is defined and produces a
normalized word.
\end{lemma}

\begin{proof}
Put $a=x_{s-1}$, let $R=(x_0,\ldots,x_{s-2})$, and let $m=\max R$.
Suppose no prefix entry lies in
\[
  [\max(0,a+1-\tau),a].
\]
Then $a\ge\tau$.  Moreover every prefix entry is at most $a$: the first
entry greater than $a$, if one existed, would have an immediate predecessor
in the forbidden interval.  Hence $m\le a-\tau$.  The final Dyck inequality
now gives
\[
  a\le x_{s-2}+\tau\le m+\tau\le a,
\]
so
\begin{equation}
  a=m+\tau,
  \qquad x_{s-2}=m.
  \label{eq:max-final-jump}
\end{equation}

By Lemma~\ref{lem:propagation}, $D$ begins with two zeros.  The prefix lambda
sum is at least $m$ by Lemma~\ref{lem:record-payment}.  Every prefix entry is
at most $a-\tau$, so every prefix-to-final kernel equals $\tau$.  Therefore
\[
  \lambda_{s-1}(D)=\tau(s-1)-a=\tau(s-2)-m.
\]
It follows that
\[
  \defc_\tau(D)\ge m+\tau(s-2)-m=\tau(s-2),
\]
contrary to~\eqref{eq:up-hyp}.

Thus the anchor exists.  Inserting $a+1$ after the first anchor gives a valid
new left adjacency.  The old Dyck inequality gives the new right adjacency,
and all other adjacencies are unchanged.
\end{proof}

\begin{lemma}[A nonfinal first extraction]
\label{lem:up-first-position}
If $D$ is not full, then it has an extractable occurrence outside its final
position.  In particular its leftmost extractable occurrence is nonfinal.
\end{lemma}

\begin{proof}
Suppose the only extractable occurrence is final.  If $x_1=0$, the first
part of Lemma~\ref{lem:boundary-ledger} gives
\[
  \defc_\tau(D)\ge\tau(s-2),
\]
a contradiction.  If $x_1>0$, Lemma~\ref{lem:propagation} makes the whole
word rapid, and Lemma~\ref{lem:rapid} gives
$\area(D)>\dinv_\tau(D)$, contradicting~\eqref{eq:strict-lower-gap}.
\end{proof}

\begin{lemma}[A second extraction after East3 failure]
\label{lem:up-second-position}
Let $e_1$ be the leftmost extractable occurrence of $D$, let $C_1$ be the
normalized word obtained by deleting it, and put
\[
  \Sigma_1=C_1:(e_1-1).
\]
If $\East_3$ fails on the final three entries of $\Sigma_1$, then $C_1$ has
an extractable occurrence outside its final two positions.
\end{lemma}

\begin{proof}
Put $e=e_1$, $p=e-1$, $C=C_1$, and $\Sigma=C:p$.  Write $c=C[-1]$.
East3 failure says
\begin{equation}
  c>p+\tau.
  \label{eq:east3-fail}
\end{equation}
Suppose $C$ has no extraction outside its final two positions.

If $C[1]=0$, the third part of Lemma~\ref{lem:boundary-ledger} and
Lemma~\ref{lem:transport} give
\[
  \defc_\tau(D)=\defc_\tau(\Sigma)\ge\tau(s-2),
\]
a contradiction.  Hence $C[1]>0$.

Lemma~\ref{lem:propagation} makes
\[
  R=C[0:-1]=(r_0,\ldots,r_{m-1}),
  \qquad m=s-2,
\]
a rapid chain.  Write $B=r_{m-1}$.  The pair $(c,p)$ has zero dinv by
\eqref{eq:east3-fail}.  Equations~\eqref{eq:rapid-dinv} and
\eqref{eq:rapid-target}, applied to the two suffix targets $c,p$, give
\[
  \dinv_\tau(\Sigma)
  \le\tau(m-1)-B+4\tau.
\]
Since $B\ge1$, transport gives
\begin{equation}
  \dinv_\tau(D)\le\tau(s+1)-2.
  \label{eq:up3-upper}
\end{equation}
The defect cutoff and~\eqref{eq:strict-lower-gap} imply
\begin{equation}
  2\dinv_\tau(D)
  \ge M_s-\tau(s-2)+3.
  \label{eq:up3-lower}
\end{equation}
For $s\ge7$, the two bounds are incompatible because
\[
 \bigl(M_s-\tau(s-2)+3\bigr)
 -2\bigl(\tau(s+1)-2\bigr)
 =\tau\frac{s(s-7)}2+7>0.
\]

It remains to treat $s=6$.  Write
\[
  R=(0,a,b,c),
  \qquad \Sigma=(0,a,b,c,z,p).
\]
Rapid growth, normalization, and East3 failure give
\begin{equation}
\begin{gathered}
  1\le a\le\tau,
  \qquad \tau<b\le a+\tau,
  \qquad a+\tau<c\le b+\tau,\\
  z\le c+\tau,
  \qquad z>p+\tau.
\end{gathered}
\label{eq:six-constraints}
\end{equation}
Put
\[
 I=\sum_{r\in\{0,a,b,c\}}
 \bigl(d_\tau(r,z)+d_\tau(r,p)\bigr),
 \qquad X=a+b+z+p.
\]
Then
\begin{equation}
  \area(\Sigma)=X+c,
  \qquad \dinv_\tau(\Sigma)=3\tau-c+I.
  \label{eq:six-stats}
\end{equation}
We use the elementary four-record estimate
\begin{equation}
  \dinv_\tau(\Sigma)-\area(\Sigma)\ge4
  \quad\Longrightarrow\quad
  X+I\le8\tau.
  \label{eq:four-record}
\end{equation}
For completeness, substitute the two triangular formulas for $d_\tau$ and
split according to the positions of $p$ and $z$.  The possible cases are
\[
\begin{array}{c|c|c}
 p<a & z\le b & X+I\le8\tau\\
 p<a & b<z\le c & X+I\le8\tau\\
 p<a & z>c & X+I\le8\tau\\
 a\le p<b & b<z\le c & X+I\le8\tau\\
 a\le p<b & z>c & X+I\le8\tau\\
 b\le p<c & z>c & X+I\le8\tau.
\end{array}
\]
The remaining cases are excluded by~\eqref{eq:six-constraints}.  Within each
row, the further subdivisions at kernel endpoints give the same affine
inequality.  The midpoint hypothesis is used in the equivalent form
\[
  I\ge X+2c-3\tau+4.
\]
This proves~\eqref{eq:four-record}.

By transport and~\eqref{eq:strict-lower-gap}, the hypothesis of
\eqref{eq:four-record} holds.  Hence
\[
 \area(\Sigma)+\dinv_\tau(\Sigma)
 =X+3\tau+I\le11\tau.
\]
This sum equals $\area(D)+\dinv_\tau(D)$, whereas the defect cutoff for
$s=6$ requires it to be at least $11\tau+1$.  This contradiction completes
the proof.
\end{proof}

\begin{remark}
The restriction $s\ge6$ is necessary for this position lemma.  For
$s=5$ and $\tau\ge7$,
\[
  D=(0,1,\tau,\tau+1,\tau+1)
\]
has defect $3\tau-1$, lies strictly below the middle, and after its first
extraction East3 fails while the only next extraction is penultimate.
\end{remark}

\section{The three DOWN lemmas}

Throughout this section, let $D$ be normalized and suppose
\begin{equation}
  0\le d=\defc_\tau(D)\le\tau(s-2)-1,
  \qquad \area(D)\le L_d.
  \label{eq:down-hyp}
\end{equation}
Assume $D$ is not a special-skeleton root.  Let $e_1$ be its leftmost
extractable occurrence and let $C_1$ be the normalized remainder.  Skeleton
and extraction tests are always made on the current prefix alone; lowered
entries already moved to the right are not included.

\begin{lemma}[The DOWN skeleton branch]
\label{lem:down-skeleton}
Under the preceding hypotheses,
\[
  C_1\text{ is full}
  \quad\Longleftrightarrow\quad
  C_1:(e_1-1)\text{ is a normalized full skeleton}.
\]
In particular, if $C_1$ is full, then
\[
  e_1-1\le C_1[-1]+\tau.
\]
\end{lemma}

\begin{proof}
Put $p=e_1-1$ and $P=C_1:p$.  If $P$ is full, then $C_1$ is full.  Indeed,
every nonfinal extraction of $C_1$ persists after appending $p$, while if
the only extraction of $C_1$ is final and the append blocks it,
Lemma~\ref{lem:propagation} makes $p$ extractable at the new right boundary.

Conversely, suppose $C_1$ is full.  It begins with $(0,0)$.  If $P$ were
invalid or not full, the second part of Lemma~\ref{lem:boundary-ledger} and
transport would give
\[
  \defc_\tau(D)=\defc_\tau(P)\ge\tau(s-2),
\]
contrary to~\eqref{eq:down-hyp}.  Hence $P$ is normalized and full, including
the asserted final adjacency.
\end{proof}

\begin{lemma}[Existence of the second DOWN extraction]
\label{lem:down-second}
If $C_1$ is not full, then it has an extractable occurrence.  Thus the second
leftmost extraction is defined on $C_1$, with no position restriction.
\end{lemma}

\begin{proof}
This is the definition of a full skeleton.
\end{proof}

Let $e_2$ be this second extracted value, let $C_2$ be the normalized
remainder, and form
\[
  T_1=C_2:(e_1-1):(e_2-1).
\]

\begin{lemma}[A nonfinal extraction after West3 failure]
\label{lem:down-third-position}
If $\West_3$ fails on the final three entries of $T_1$, then $C_2$ has an
extractable occurrence outside its final position.
\end{lemma}

\begin{proof}
Put
\[
  p=e_1-1,
  \qquad q=e_2-1,
  \qquad c=C_2[-1].
\]
West3 failure is
\begin{equation}
  p>c+\tau.
  \label{eq:west3-fail}
\end{equation}
Suppose $C_2$ has no extraction outside its final position.  Two applications
of transport and the weak lower-half inequality give
\begin{equation}
  \dinv_\tau(T_1)-\area(T_1)
  =\dinv_\tau(D)-\area(D)+4\ge4.
  \label{eq:down-gap}
\end{equation}

If $C_2[1]=0$, the fourth part of Lemma~\ref{lem:boundary-ledger} gives
\[
  \defc_\tau(D)=\defc_\tau(T_1)\ge\tau(s-2),
\]
a contradiction.  Hence $C_2[1]>0$.

Lemma~\ref{lem:propagation} makes
\[
  C_2=R=(r_0,\ldots,r_{m-1}),
  \qquad m=s-2,
  \qquad c=r_{m-1},
\]
a rapid chain.  The unique predecessor of $e_1=p+1$ has value at least
$p+1-\tau>c$.  Since every entry remaining in $C_2$ is at most $c$, this
predecessor must be the second extracted value $e_2$.  Therefore
\begin{equation}
  q=e_2-1\ge p-\tau>c.
  \label{eq:q-above-c}
\end{equation}

No pair from $R$ to $p$ contributes to dinv.  The rapid estimates give
\[
 \dinv_\tau(T_1)
 \le\tau(m-1)-c+2\tau+\tau
 =\tau(m+2)-c.
\]
By integrality in~\eqref{eq:west3-fail} and~\eqref{eq:q-above-c},
\[
 \area(T_1)=\area(R)+p+q
 \ge\area(R)+2c+\tau+2.
\]
Equation~\eqref{eq:short-rapid} now gives
\[
 \area(T_1)-\dinv_\tau(T_1)
 \ge\area(R)+3c+2-\tau(m+1)>0,
\]
contradicting~\eqref{eq:down-gap}.
\end{proof}

\section{Definition and totality of the maps}

We now define the maps only in the defect range under consideration.  The
two exceptional global words have defect $\tau(s-2)$ and do not occur.

\subsection{The UP map}

Let $D$ be strictly below the middle.
\begin{enumerate}
\item If $D=C:a$ is full, set
\[
  \UP(D)=\inj_{a+1}(C).
\]
\item Otherwise extract $e_1$, leaving $C_1$, and form
$\Sigma_1=C_1:(e_1-1)$.  If East3 succeeds, apply it to the final window,
raise the final two local output entries by one, and inject them from right
to left into the remaining prefix.
\item If East3 fails, extract $e_2$ from $C_1$, leaving $C_2$, and form
\[
  \Sigma_2=C_2:(e_1-1):(e_2-1).
\]
Apply East5 to the final window, raise its final three output entries by one,
and inject them from right to left into the remaining prefix.
\end{enumerate}

Lemmas~\ref{lem:up-skeleton}--\ref{lem:up-second-position} give all required
anchors and position bounds.  Assumption~\ref{ass:east5} supplies East5
nonfailure.  A direct check of the local inequalities shows that the raised
outputs have prefix anchors and that the resulting word is normalized.
Thus $\UP$ is defined on every word satisfying~\eqref{eq:up-hyp}.

\subsection{The DOWN map}

Let $D$ be weakly below the middle and not a special skeleton.
\begin{enumerate}
\item Extract $e_1$, leaving $C_1$.  If $C_1$ is full, set
\[
  \DOWN(D)=C_1:(e_1-1).
\]
\item Otherwise extract $e_2$, leaving $C_2$, and form
\[
  T_1=C_2:(e_1-1):(e_2-1).
\]
If West3 succeeds, apply it and raise and inject the final local output entry.
\item If West3 fails, extract $e_3$ from $C_2$, leaving $C_3$, and form
\[
  T_2=C_3:(e_1-1):(e_2-1):(e_3-1).
\]
Apply West5, raise its final two output entries by one, and inject them from
right to left.
\end{enumerate}

Lemmas~\ref{lem:down-skeleton}--\ref{lem:down-third-position} give the
required skeleton, extraction, and position statements.
Assumption~\ref{ass:west5} supplies West5 nonfailure.  The local inequalities
give the remaining anchor and normalization checks.  Thus $\DOWN$ is defined
on every weakly lower-half nonskeleton in the stated defect range.

\section{Statistics and invertibility}

\begin{lemma}[Local preservation]
\label{lem:local-preservation}
East3 preserves area and dinv.  On its relevant domain, East5 also preserves
area and dinv.
\end{lemma}

\begin{proof}
East3 is the identity.  East5 only rearranges five entries, so it preserves
area.  Every interaction with an earlier prefix occurrence is unchanged
because the multiset of the five entries is unchanged.  It remains to compare
the six internal pairs among the four moved entries.  Substitution of the
two formulas defining $d_\tau$ in the two East5 cases gives
\[
  \sum_{i<j}d_\tau(w_i,w_j)
  =\sum_{i<j}d_\tau(w'_i,w'_j),
\]
where $w,w'$ are the input and output windows.  In the second case,
$\bk_\tau$ fixes or switches the middle pair according to precisely the gap
at which the two formulas exchange.  This proves local dinv preservation.
\end{proof}

\begin{proposition}[Statistic change]
\label{prop:statistics}
Whenever $\UP$ is defined,
\[
 \area(\UP(D))=\area(D)+1,
 \qquad
 \dinv_\tau(\UP(D))=\dinv_\tau(D)-1.
\]
Consequently $\UP$ preserves defect.  Once invertibility is established,
$\DOWN$ has the reverse statistic change.
\end{proposition}

\begin{proof}
If a branch makes $k$ extractions before its local East move, repeated use of
Lemma~\ref{lem:transport} gives stage statistics
\[
  (\area(D)-k,\dinv_\tau(D)+k).
\]
The local East move preserves them.  Extraction and injection are inverse:
raising a terminal $p$ to $p+1$ and injecting it reverses one extraction
transport and therefore changes the statistics by $(+1,-1)$.  UP makes
$k+1$ such raised injections.  Its net change is consequently $(+1,-1)$.
The full-skeleton branch is the case $k=0$.
\end{proof}

We next establish invertibility.  We use the following two elementary facts.
First, entries injected right to left in an admissible local configuration
are successively extracted in reverse injection order.  Second, successive
leftmost extracted values satisfy
\begin{equation}
  e_1\le e_2+\tau.
  \label{eq:successive-extraction}
\end{equation}
Indeed, if the reverse inequality held, deleting $e_1$ would not improve
either the predecessor count or splice of $e_2$, unless the two were
adjacent; adjacency gives the original Dyck inequality and again contradicts
the reverse inequality.

The local maps have the corresponding recognition property.  If
\begin{equation}
  \East_5(*,b,c,x,y)=(*,r,j,k,y),
  \label{eq:east5-symbolic}
\end{equation}
then
\begin{equation}
  j>r+\tau,
  \qquad
  \West_5(*,r,j,k,y)=(*,b,c,x,y).
  \label{eq:east-west-inverse}
\end{equation}
In East5 case 2b, $(r,j,k)=(x,c,b)$ and the first inequality is the original
East3 failure.  In case 2a it follows from the defining gap and the possible
$\bk_\tau$ switch.  The second identity, and its converse, follow by the same
two-case check.

\begin{proposition}[Invertibility]
\label{prop:inverse}
On the stated domains,
\[
  \DOWN\circ\UP=\mathrm{id},
  \qquad
  \UP\circ\DOWN=\mathrm{id}.
\]
\end{proposition}

\begin{proof}
We pair the three UP branches with the three DOWN branches.

\smallskip
\noindent\emph{UP skeleton, then DOWN.}
Write the full skeleton as $S=C:a$.  The prefix $C$ is full: an extraction
in $C$ would persist in $S$, or its obstruction would propagate to an
extraction of $a$.  UP injects $a+1$ into $C$.  The inserted occurrence is
the first extractable one.  DOWN extracts it, recognizes the full prefix
$C$, and appends $(a+1)-1=a$.

\smallskip
\noindent\emph{UP East3, then DOWN West3.}
Write $D=P:b$ and $P=\inj_{x+1}(P')$.  The East3 stage is $P':b:x$, and the
UP output is
\[
  \inj_{b+1}\bigl(\inj_{x+1}(P')\bigr).
\]
DOWN extracts $b+1$ and then $x+1$, so its skeleton branch does not apply.
The word $P':b$ is normalized because the original extraction was valid;
therefore the West3 condition holds.  West3 is the identity, and injecting
$x+1$ into $P':b$ recovers $P:b$.

\smallskip
\noindent\emph{UP East5, then DOWN West5.}
After two UP extractions, write the stage as $P':b:c:x:y$, and use
\eqref{eq:east5-symbolic}.  The raised entries $y+1,k+1,j+1$ inject before
the retained $r$, so the UP output is $Q:r$.  DOWN successively extracts
$j+1,k+1,y+1$.  After the second extraction its West3 window is $(r,j,k)$,
which fails by~\eqref{eq:east-west-inverse}.  The third extraction reaches
$P':r:j:k:y$, West5 returns $P':b:c:x:y$, and injecting $y+1$ and then
$x+1$ reverses the original two extractions.

\smallskip
For the other composition, first suppose DOWN takes its skeleton branch.
It extracts $x+1$, leaves a full prefix $C$, and returns $C:x$.
Lemma~\ref{lem:down-skeleton} says $C:x$ is full, so UP removes $x$, raises
it, and injects $x+1$ back into $C$.

If DOWN takes West3 after extracting $x+1,y+1$, write its stage as
$P':b:x:y$.  It injects $y+1$ before the retained $x$, producing $Q:x$.
UP extracts $y+1$, so its skeleton branch does not apply.  Inequality
\eqref{eq:successive-extraction} gives $x\le y+\tau$, exactly the East3
condition.  East3 is the identity, and the two reverse injections recover
the DOWN input.

Finally suppose DOWN takes West5 after extracting $x+1,y+1,z+1$.  Write
\[
  \West_5(*,b,x,y,z)=(*,j,k,r,z).
\]
DOWN injects $z+1,r+1$ before the retained $j,k$, producing $Q:j:k$.
UP extracts $r+1$.  The East3 window $(j,k,r)$ fails because $k>r+\tau$,
the converse recognition inequality to~\eqref{eq:east-west-inverse}.  UP
extracts $z+1$, East5 returns $(*,b,x,y,z)$, and injecting
$z+1,y+1,x+1$ reverses the three original extractions.
\end{proof}

\begin{corollary}
Whenever $\DOWN$ is defined,
\[
 \area(\DOWN(D))=\area(D)-1,
 \qquad
 \dinv_\tau(\DOWN(D))=\dinv_\tau(D)+1,
\]
and $\DOWN$ preserves defect.
\end{corollary}

\begin{proof}
This follows from Proposition~\ref{prop:statistics} and
Proposition~\ref{prop:inverse}.
\end{proof}

\section{Proof of the decomposition formula}

\begin{proof}[Proof of Theorem~\ref{thm:main}]
Put $B=\tau(s-2)-1$.  We first locate the roots.  If $S$ is a special
skeleton of defect $d\le B$, Assumption~\ref{ass:skeleton-area} gives
\[
  2\area(S)+d\le3d\le3B.
\]
For $s\ge6$,
\[
  M_s-3B
  =\tau\frac{(s-3)(s-4)}2+3>0.
\]
Thus $2\area(S)+d<M_s$, so every root is strictly below the middle.  Since
the exceptional skeleton has defect $B+1$, every full skeleton in the range
is special.

Fix a defect $d\le B$.  The UP lemmas, Assumption~\ref{ass:east5}, and the
local validity checks show that UP is defined on every word of defect $d$
with area at most $L_d-1$.  The DOWN lemmas, Assumption~\ref{ass:west5}, and
the local validity checks show that DOWN is defined on every weakly
lower-half nonskeleton of defect $d$.  By Propositions~\ref{prop:statistics}
and~\ref{prop:inverse}, they preserve defect, move area by one in opposite
directions, and are inverse.

Starting from any weakly lower-half word, iterate DOWN until it is undefined.
Area decreases at every step, so the process terminates.  It can terminate
only at a full, hence special, skeleton.  Invertibility makes this root and
the entire descent unique.  Conversely, starting at a special skeleton $S$
and iterating UP produces exactly one word at each area
\[
  \area(S),\area(S)+1,\ldots,L_d.
\]
Thus the weak lower half of the defect-$d$ layer is the disjoint union of
these rooted strings.

Put $N=M_s-d$.  Along the string rooted at $S$, the word of area $j$ has
dinv $N-j$.  Its lower-half contribution is
\[
  \sum_{j=\area(S)}^{\lfloor N/2\rfloor}q^{N-j}t^j.
\]
The assumed $q,t$ symmetry supplies the reflected upper half, with the
central term counted once when $N$ is even.  Since
$\dinv_\tau(S)=N-\area(S)$, the complete string contributes
\begin{align*}
  \sum_{j=\area(S)}^{\dinv_\tau(S)}q^{N-j}t^j
  &=
  \frac{
  q^{\dinv_\tau(S)+1}t^{\area(S)}
  -q^{\area(S)}t^{\dinv_\tau(S)+1}
  }{q-t}.
\end{align*}
Summing over all special skeletons of defect at most $B$ proves the theorem.
\end{proof}

\end{document}
